% !TeX program = tectonic
\documentclass[11pt,a4paper]{article}
\usepackage{fontspec}
\usepackage[russian]{babel}
\babelfont{rm}[Language=Default]{Liberation Serif}
\babelfont[Language=Default]{sf}{Carlito}
\babelfont[Language=Default]{tt}{DejaVu Sans Mono}
\usepackage{amsmath,amssymb,amsthm}
\usepackage{geometry}
\geometry{a4paper, top=2.4cm, bottom=2.4cm, left=2.4cm, right=2.4cm}
\usepackage{booktabs,tabularx,enumitem,xcolor,tcolorbox}
\tcbuselibrary{breakable,skins}
\definecolor{linkblue}{HTML}{1F4E79}
\definecolor{statgreen}{HTML}{2F6B3A}
\definecolor{goldacc}{HTML}{8B7E5A}
\usepackage[colorlinks=true,linkcolor=linkblue,urlcolor=linkblue,unicode=true]{hyperref}
\theoremstyle{plain}
\newtheorem{theorem}{Теорема}
\newtheorem{lemma}[theorem]{Лемма}
\newtheorem{proposition}[theorem]{Предложение}
\newtheorem{corollary}[theorem]{Следствие}
\theoremstyle{definition}
\newtheorem{definition}[theorem]{Определение}
\theoremstyle{remark}
\newtheorem{remark}[theorem]{Замечание}
\newtcolorbox{statusbox}[1][]{colback=statgreen!5,colframe=statgreen!75,
  title={\textbf{Сводка доказанного:} #1},fonttitle=\small,boxrule=0.7pt,arc=2pt,breakable}
\newtcolorbox{databox}[1][]{colback=goldacc!6,colframe=goldacc!80,
  title={\textbf{Данные протокола:} #1},fonttitle=\small,boxrule=0.7pt,arc=2pt,breakable}
\newtcolorbox{verifybox}[1][]{colback=linkblue!5,colframe=linkblue!80,
  title={\textbf{Верификация:} #1},fonttitle=\small,boxrule=0.7pt,arc=2pt,breakable}
\widowpenalty=10000
\clubpenalty=10000
\setlength{\parskip}{0.35em}
\newcommand{\CC}{\mathbb{C}}
\newcommand{\RR}{\mathbb{R}}
\newcommand{\ZZ}{\mathbb{Z}}
\newcommand{\QQ}{\mathbb{Q}}
\newcommand{\Ga}{\Gamma}
\newcommand{\Bfun}{\mathrm{B}}
\newcommand{\im}{\operatorname{Im}}
\newcommand{\re}{\operatorname{Re}}
\newcommand{\lcm}{\operatorname{lcm}}
\newcommand{\SNF}{\operatorname{SNF}}
\newcommand{\Jac}{\operatorname{Jac}}
\newcommand{\rank}{\operatorname{rank}}
\newcommand{\Ind}{\operatorname{Index}}
\newcommand{\disc}{\operatorname{disc}}
\begin{document}
\begin{center}
{\LARGE\bfseries Сертификат F: рационализация}\\[0.4em]
{\large априорная граница знаменателей и точная LLL}\\[0.6em]
{\normalsize Исаев Исхак Хамзатович}\\[0.2em]
{\small Программа <<Динамический принцип>> --- лаборатория \texttt{hodge-laboratory}}\\[0.15em]
{\small Отдельная монография теоремы · Монография 14/16}
\end{center}
\vspace{0.4em}
{\color{linkblue}\hrule height 0.8pt}
\vspace{0.8em}

\section{Постановка}

Сертификат F решает задачу рационализации: по численным измерениям
восстановить точные рациональные значения с априорной гарантией
единственности. Ключ --- граница знаменателей, вычисляемая до
измерений.

\begin{theorem}[F1--F4]\label{th:Fstand}
\begin{enumerate}[leftmargin=2em]
\item[(F1)] Единственность: если $2\varepsilon<1/Q^2$, где $Q$ ---
граница знаменателей, то в интервале ширины $2\varepsilon$ лежит не
более одного рационального числа со знаменателем $\le Q$. При $100$
цифрах сертифицированная граница $Q<7.07\cdot10^{49}$.
\item[(F2)] Априорная граница вычисляется из геометрии: K3 ---
$4|\det G_8|=256$ (правило Крамера), тор --- $1$, Клейн --- $2$
(в $\tau$-базисе) и $1$ (в целых).
\item[(F3)] Целочисленный $q$-скан восстанавливает значения без
плавающей точки: сложность $O(nQ)$, критерий приёма $d_q\le q/2$.
\item[(F4)] Точная LLL: $\im\tau=\sqrt7/2$ восстановлено с
минимальным многочленом $4X^2-7$; полиномы степени $2$
единственны.
\end{enumerate}
\end{theorem}

\begin{proof}
(F1) Диофантова оценка: два несократимых $\frac pq\ne\frac{p'}{q'}$
со знаменателями $\le Q$ различаются не менее чем на
$\frac1{qQ'}\ge\frac1{Q^2}$; интервал ширины $<1/Q^2$ не может
содержать оба.
(F2) По правилу Крамера координаты целевого класса --- отношения
определителей матриц из точной матрицы пересечений; знаменатели
делят $\det G_8=64$; граница $4\cdot64=256$ с запасом.
(F3) Для каждого кандидата $q\le Q$: $d_q=\min_n|qx-n|$ в целых
умножениях; приём при $2d_q\le1$.
(F4) LLL на решётке периода Клейна: точные операции в
$\QQ(\sqrt{-7})$ дают $\im\tau=\sqrt7/2$; минимальный многочлен
$4X^2-7$ единственен; $j(\tau)=-3375$ с отклонением
$1.99\cdot10^{-117}$.
\end{proof}

\begin{databox}[сквозной прогон K3]
$49$ измерений $\to$ восстановление $[Z]=[target]$ точно;
максимальный знаменатель $4$; априорная граница $Q=256$ (запас
$\times64$); тор: степени $(2,2)$; Клейн: $\re\tau=1/2$,
$\im\tau=\sqrt7/2$, $\mathrm{tr}\,t=-1$, $\mathrm{norm}\,t=2$,
$j=-3375$; остаток $\im\tau$: $7.75\cdot10^{-121}$.
Отрицательный тест: значение $4\pi^2$ \emph{отвергнуто} сканом ---
согласовано с трансцендентностью по Линдеману.
\end{databox}

\begin{remark}[роль отрицательного теста]
Отвержение $4\pi^2$ --- контрольная проверка мощности метода: скан
не принимает <<похожие>> значения. Согласие с теоремой Линдемана
(трансцендентность $\pi$) --- независимое подтверждение
корректности критерия приёма.
\end{remark}

\begin{statusbox}[итог]
Сертификат F принят ($5/5$): рационализация с гарантированной
границей работает end-to-end; отрицательный тест пройден.
\end{statusbox}

\begin{verifybox}[как воспроизвести]
\texttt{python3 laboratory.py} $\to$ <<Сертификаты $\to$ F>>;
конструктор: произвольное значение и граница $Q$ --- лаборатория
восстанавливает и печатает вердикт.
\end{verifybox}

\vfill
{\color{linkblue}\hrule height 0.6pt}
\vspace{0.5em}
{\small\color{gray}
Лицензия: индивидуальная исключительная лицензия Исаева Исхака Хамзатовича.
Все права на вычисления, формулы и тексты принадлежат автору.
Верификация: \texttt{python3 laboratory.py} (пункт <<Стенды>>/<<Сертификаты>>),
Lean: \texttt{verification/lean/HodgeLaboratory.lean}.
}
\end{document}
